
\documentclass{article}
\usepackage{colortbl}
\usepackage{makecell}
\usepackage{multirow}
\usepackage{supertabular}

\begin{document}

\newcounter{utterance}

\centering \large Interaction Transcript for game `hot\_air\_balloon', experiment `air\_balloon\_survival\_en\_negotiation\_easy', episode 2 with Qwen3{-}32B{-}t0.0.
\vspace{24pt}

{ \footnotesize  \setcounter{utterance}{1}
\setlength{\tabcolsep}{0pt}
\begin{supertabular}{c@{$\;$}|p{.15\linewidth}@{}p{.15\linewidth}p{.15\linewidth}p{.15\linewidth}p{.15\linewidth}p{.15\linewidth}}
    \# & $\;$A & \multicolumn{4}{c}{Game Master} & $\;\:$B\\
    \hline

    \theutterance \stepcounter{utterance}  
    & & \multicolumn{4}{p{0.6\linewidth}}{
        \cellcolor[rgb]{0.9,0.9,0.9}{
            \makecell[{{p{\linewidth}}}]{
                \texttt{\tiny{[P1$\langle$GM]}}
                \texttt{You are participating in a collaborative negotiation game.} \\
\\ 
\texttt{Together with another participant, you must agree on a single set of items that will be kept. Each of you has your own view of how much each item matters to you (importance). You do not know how the other participant values the items. Additionally, you are given the effort each item demands.} \\
\texttt{You may only agree on a set if the total effort of the selected items does not exceed a shared limit:} \\
\\ 
\texttt{LIMIT = 3288} \\
\\ 
\texttt{Here are the individual item effort values:} \\
\\ 
\texttt{Item effort = \{"C08": 329, "C55": 41, "A43": 417, "A67": 75, "B42": 390, "A62": 887, "B19": 808, "C43": 151, "C73": 849, "C10": 129, "B36": 350, "A18": 118, "C61": 630, "C67": 602, "C42": 801\}} \\
\\ 
\texttt{Here is your personal view on the importance of each item:} \\
\\ 
\texttt{Item importance values = \{"C08": 138, "C55": 583, "A43": 868, "A67": 822, "B42": 783, "A62": 65, "B19": 262, "C43": 121, "C73": 508, "C10": 780, "B36": 461, "A18": 484, "C61": 668, "C67": 389, "C42": 808\}} \\
\\ 
\texttt{Goal:} \\
\\ 
\texttt{Your goal is to negotiate a shared set of items that benefits you as much as possible (i.e., maximizes total importance to YOU), while staying within the LIMIT. You are not required to make a PROPOSAL in every message {-} you can simply negotiate as well. All tactics are allowed!} \\
\\ 
\texttt{Interaction Protocol:} \\
\\ 
\texttt{You may only use the following structured formats in a message:} \\
\\ 
\texttt{PROPOSAL: \{'A', 'B', 'C', …\}} \\
\texttt{Propose keeping exactly those items.} \\
\texttt{REFUSE: \{'A', 'B', 'C', …\}} \\
\texttt{Explicitly reject opponent's proposal.} \\
\texttt{ARGUMENT: \{'...'\}} \\
\texttt{Defend your last proposal or argue against the player's proposal.} \\
\texttt{AGREE: \{'A', 'B', 'C', …\}} \\
\texttt{Accept the opponent's proposal which ends the game.} \\
\texttt{STRATEGIC REASONING: \{'...'\}} \\
\texttt{	Describe your strategic reasoning or anticipation explaining your choice of action. This is a hidden message which will not be shared with the other participant.} \\
\\ 
\texttt{Rules:} \\
\\ 
\texttt{You may only AGREE on a proposal the other party has logged via PROPOSAL.} \\
\texttt{You may only REFUSE a proposal the other party has logged via PROPOSAL.} \\
\texttt{Total effort of any PROPOSAL or AGREE set must be ≤ LIMIT.} \\
\texttt{Do NOT reveal your hidden importance scores.} \\
\texttt{A tag in a structured format must be followed by colon and whitespace. The argument must be a python set containing 0 or more strings.} \\
\texttt{So, it must be of the form TAG: \{...\}} \\
\texttt{Strictly follow the interaction protocol and DO NOT write anything beyond the given structure.} \\
\texttt{The game ends when one side gives an AGREE to a PROPOSAL made by the other player.} \\
\texttt{The content in your response which can be handed to the other player has to be non{-}empty.} \\
\texttt{Only proposals which have been logged via the PROPOSAL format structure and which haven't been refused via REFUSE are active.} \\
\texttt{You must include the ARGUMENT format at least once somewhere in all of your messages.} \\
\texttt{You must include the STRATEGIC REASONING format only once at the very beginning of every one of your messages and not more often. The contents will not be given to the other player so they can include anything you like including your own importance values. Here you should reason step by step to come up with you next move.} \\
            }
        }
    }
    & \\ \\

    \theutterance \stepcounter{utterance}  
    & \multicolumn{4}{p{0.6\linewidth}}{
        \cellcolor[rgb]{0.9,0.9,0.9}{
            \makecell[{{p{\linewidth}}}]{
                \texttt{\tiny{[P1$\rangle$GM]}}
                \texttt{<think>} \\
\texttt{Okay, let's see. I need to figure out the best set of items to propose that maximizes my total importance without exceeding the effort limit of 3288. First, I should calculate the total effort for different combinations. Let me list out the items with their effort and importance values.} \\
\\ 
\texttt{Looking at the importance values, the highest ones are A43 (868), C10 (780), B42 (783), A18 (484), C61 (668), C08 (138), etc. But I have to balance that with the effort each item takes. For example, A43 has high importance but also high effort (417). C10 has high importance (780) and moderate effort (129). B42 is 783 importance and 390 effort. Hmm, that's a lot. Let me check the total effort if I take some of these.} \\
\\ 
\texttt{Let me start by considering the items with the highest importance per unit of effort. Maybe that's a good way to maximize the total importance without going over the limit. Let's calculate the ratio of importance to effort for each item:} \\
\\ 
\texttt{{-} A43: 868 / 417 ≈ 2.08} \\
\texttt{{-} C10: 780 / 129 ≈ 6.05} \\
\texttt{{-} B42: 783 / 390 ≈ 2.01} \\
\texttt{{-} A18: 484 / 118 ≈ 4.10} \\
\texttt{{-} C61: 668 / 630 ≈ 1.06} \\
\texttt{{-} C08: 138 / 329 ≈ 0.42} \\
\texttt{{-} A67: 822 / 75 ≈ 10.96} \\
\texttt{{-} C55: 583 / 41 ≈ 14.22} \\
\texttt{{-} B36: 461 / 350 ≈ 1.32} \\
\texttt{{-} C43: 121 / 151 ≈ 0.80} \\
\texttt{{-} C73: 508 / 849 ≈ 0.60} \\
\texttt{{-} B1} \\
            }
        }
    }
    & & \\ \\

    \theutterance \stepcounter{utterance}  
    & & \multicolumn{4}{p{0.6\linewidth}}{
        \cellcolor[rgb]{0.9,0.9,0.9}{
            \makecell[{{p{\linewidth}}}]{
                \texttt{\tiny{[P1$\langle$GM]}}
                \texttt{Your response did not start with the proper strategic reasoning tag at the very beginning of your response. The very first structured format must be of the form STRATEGIC REASONING: \{...\}. Try again.} \\
            }
        }
    }
    & \\ \\

    \theutterance \stepcounter{utterance}  
    & \multicolumn{4}{p{0.6\linewidth}}{
        \cellcolor[rgb]{0.9,0.9,0.9}{
            \makecell[{{p{\linewidth}}}]{
                \texttt{\tiny{[P1$\rangle$GM]}}
                \texttt{<think>} \\
\texttt{Okay, I need to make sure I start with the STRATEGIC REASONING tag. Let me think through this again. My goal is to maximize my total importance while staying under the effort limit. The items with the highest importance per unit effort are A67 (10.96), C55 (14.22), and C10 (6.05). But I also need to consider the total effort. Let me check the effort of these top items. A67 is 75, C55 is 41, and C10 is 129. Adding them up: 75 + 41 + 129 = 245. That's way under 3288. Maybe I can add more high{-}value items. Next, A43 has high importance (868) but effort 417. Adding that would bring total to 245 + 417 = 662. Still under. Then B42 is 783 importance with 390 effort. Adding that would make 662 + 390 = 1052. Still okay. Then A18 is 484 with 118 effort. Total becomes 1052 + 118 = 1170. Next, C61 is 668 with 630 effort. Adding that gives 1170 + 630 = 1800. Still under. Then C08 is 138 with 329. Total 1800 + 329 = 2129. Then B36 is 461 with 350. Total 2129 + 350 = 2479. Then C43 is 121 with 151. Total 2479 + 151 = 2630. Then C73 is 508 with 849. Adding that would exceed the limit (2630 + 849 = 3479 > 3288). So maybe skip C73. What about C67? 389 importance with 602 effort. 26} \\
            }
        }
    }
    & & \\ \\

    \theutterance \stepcounter{utterance}  
    & & \multicolumn{4}{p{0.6\linewidth}}{
        \cellcolor[rgb]{0.9,0.9,0.9}{
            \makecell[{{p{\linewidth}}}]{
                \texttt{\tiny{[P1$\langle$GM]}}
                \texttt{Your response did not start with the proper strategic reasoning tag at the very beginning of your response. The very first structured format must be of the form STRATEGIC REASONING: \{...\}. Try again.} \\
            }
        }
    }
    & \\ \\

    \theutterance \stepcounter{utterance}  
    & \multicolumn{4}{p{0.6\linewidth}}{
        \cellcolor[rgb]{0.9,0.9,0.9}{
            \makecell[{{p{\linewidth}}}]{
                \texttt{\tiny{[P1$\rangle$GM]}}
                \texttt{<think>} \\
\texttt{Okay, I need to make sure I start with the STRATEGIC REASONING tag. Let me outline my plan step by step. First, I'll calculate the importance per unit effort for each item to prioritize those with the highest ratios. The top items here are C55 (583/41 ≈14.22), A67 (822/75≈10.96), and C10 (780/129≈6.05). These should be included first. Next, I'll add other high{-}importance items while keeping the total effort under 3288. Let's check the cumulative effort: C55 (41) + A67 (75) + C10 (129) = 245. Then adding A43 (417) brings it to 662. B42 (390) adds up to 1052. A18 (118) makes 1170. C61 (630) totals 1800. C08 (329) adds to 2129. B36 (350) reaches 2479. C43 (151) is 2630. C67 (602) would exceed the limit (2630+602=3232, which is under). Wait, 2630 + 602 = 3232, which is under 3288. Then adding C67 (602) gives 3232. Then check if there's room for more. The next item could be C73 (849), but 3232 + 849 = 4081, which is over. So maybe stop there. The total importance would be C55 (583) + A67 (822) + C10 (780) + A43 (868) + B42 (783) + A18 (484) + C61 (668) + C08 (138) + B36 (461) + C43 (} \\
            }
        }
    }
    & & \\ \\

    \theutterance \stepcounter{utterance}  
    & & & \multicolumn{2}{p{0.3\linewidth}}{
        \cellcolor[rgb]{0.9,0.9,0.9}{
            \makecell[{{p{\linewidth}}}]{
                \texttt{\tiny{[GM$|$GM]}}
                \texttt{abort} \\
            }
        }
    }
    & & \\ \\

    \theutterance \stepcounter{utterance}  
    & & & \multicolumn{2}{p{0.3\linewidth}}{
        \cellcolor[rgb]{0.9,0.9,0.9}{
            \makecell[{{p{\linewidth}}}]{
                \texttt{\tiny{[GM$|$GM]}}
                \texttt{end game} \\
            }
        }
    }
    & & \\ \\

\end{supertabular}
}

\end{document}
